\par\medskip
\noindent
\underline{Proposición} 3.2. \underline{Sea $S$ un esquema entero
regular, localmente noetheriano, de dimensión} $1$; \underline{sea $K$ su cuerpo de
funciones racionales; sea $A_K$ un esquema abeliano sobre $K$; y sea
$A$ su modelo de Néron sobre $S$, que es un esquema en grupos liso y
separado sobre $S$} (\underline{y cuasiproyectivo sobre $S$ si $S$ es noetheriano
[23, VIII, 2, 3\textdegree]). Son equivalentes las condiciones
siguientes:}

\smallskip
\noindent
(i) \underline{Para todo $s\in S$, $A_s^o$ posee rango unipotente cero;
es decir, es una extensión de un esquema abeliano por un toro (2.1.5).}

\smallskip
\noindent
(ii) \underline{Existe un esquema en grupos $G$, liso, separado y de
tipo finito sobre $S$, que prolonga $A_K$ y cuyas fibras poseen rango
unipotente cero.}

\smallskip
\noindent
\underline{Además, si $G$ es como en (ii), el único morfismo
$G\longrightarrow A$ que induce el isomorfismo dado sobre las fibras
genéricas es un isomorfismo de $G$ sobre un subgrupo abierto de $A$;
induce, por tanto, un isomorfismo de las componentes neutras}
\[
 G^o\xrightarrow{\ \sim\ }A^o.
\]

\medskip
\noindent
3.2.1. Observemos de inmediato que la definición del modelo de Néron
$A$ de $A_K$ es la misma que en el caso local recordado en 1.1. La
existencia de este modelo de Néron, y el hecho de que sea separado y de
tipo finito sobre $S$, son bien conocidos y, de hecho, se reducen de
inmediato al caso local. Naturalmente, para todo punto cerrado $s\in S$,
\[
 S'=\operatorname{Spec}(\mathcal O_{S,s})
\]
es un rasgo, y $A\times_S S'$ es precisamente el modelo de Néron de
$A_K$ sobre $S'$ en el sentido considerado en 1.1.

Demostremos 3.2. Es trivial que (i) implica (ii), pues basta tomar
$G=A$. Queda por demostrar que, si $G$ es como en (ii), entonces
$G\longrightarrow A$ es una inmersión abierta. Este es, en efecto, un
caso particular de 3.1 e).

\medskip
\noindent
\underline{Corolario} 3.3. \underline{Con la notación de 3.2,
supongamos que las componentes neutras de las fibras del modelo de
Néron poseen rango unipotente cero. Sea $S'$ un esquema que satisface
las mismas hipótesis que $S$; sea $K'$ su cuerpo de funciones
racionales; sea $f:S'\longrightarrow S$ un morfismo dominante; sea}
\[
 A_{K'}\simeq A_K\otimes_KK'
 \simeq \text{la fibra genérica de }A\times_SS',
\]
\underline{sea $A'$ el modelo de Néron de $A_{K'}$, y considérese el
único morfismo}
\[
 u:A\times_SS'\longrightarrow A'
\]
\underline{que induce la identidad sobre las fibras genéricas. Entonces
$u$ es una inmersión abierta y, por tanto, induce un isomorfismo sobre
las componentes neutras}
\begin{equation}\tag{3.3.1}
 A^o\times_SS'\xrightarrow{\ \sim\ }{A'}^o.
\end{equation}
En efecto, basta aplicar 3.2 sobre $S'$, tomando
$G\simeq A\times_SS'$.

\medskip
\noindent
3.3.2. Así, las propiedades equivalentes de 3.2 implican una propiedad
parcial de «estabilidad» de los modelos de Néron. Conviene tener presente
que, cuando $A_K$ no posee buena reducción, es decir, cuando $A$ no es
un esquema abeliano, no es cierto que, para todo $S'$ como en 3.3
(incluso para $S'$ finito sobre $S$ y $K'/K$ separable), el propio
morfismo
\[
 A\times_SS'\longrightarrow A'
\]
sea un isomorfismo, ni que induzca isomorfismos sobre los grupos de
componentes conexas de las fibras (véase 11.10 más adelante para el
comportamiento de estos grupos de componentes conexas bajo extensión
de la base). Estas observaciones, reforzadas por 3.9 más adelante,
justifican la terminología siguiente, introducida por Mumford (y que
también parece concordar con las nociones del mismo nombre estudiadas
en su libro [17]).

\medskip
\noindent
\underline{Definición} 3.4. \underline{Con la notación de 3.2, se dice
que el esquema abeliano $A_K$ sobre $K$ posee reducción semiestable
sobre $S$ si se satisfacen las condiciones equivalentes (i) y (ii) de
3.2.}

\medskip
\noindent
3.4.0. Obsérvese que, en la forma (ii), esta condición no invoca la
noción de modelo de Néron. Además, por localización, se reduce de
inmediato al caso en que $S$ es un rasgo, como se deduce, por ejemplo,
de 3.2.1. Por otra parte, puesto que la formación del modelo de Néron
conmuta con el cambio de base de $S$ a su henselización o a su
henselización estricta, la noción se reduce incluso al caso en que $S$
es un rasgo estrictamente local.
